\section{Numerical induction before successor existence}\label{un:core}

The numerical representatives and saturated codes below are adapted from
the supplied source-only draft~\cite{draftC2}. The new dependency order is
essential: we prove induction and global minimization before constructing
successors. This makes a source-free unary matrix sufficient.

Throughout this section membership is the sole primitive binary relation.
Define
\begin{align*}
 u\Eq v&\ :\Longleftrightarrow\ \forall z\,(z\in u\leftrightarrow z\in v),\\
 \Sat(a)&\ :\Longleftrightarrow\ \forall u\forall v\,
 [u\Eq v\to(u\in a\leftrightarrow v\in a)],\\
 \Tr(a)&\ :\Longleftrightarrow\ \forall u\forall v\,
 [(u\in v\land v\in a)\to u\in a].
\end{align*}
The relation $\Eq$ is an equivalence relation by logic alone. It is always
a congruence in the second coordinate of membership. In the first
coordinate it is a congruence when the containing object is saturated.
We never assume that it is a congruence throughout the ambient universe.

Let $\operatorname{Lin}(a)$ say that membership among the members of $a$
is irreflexive and transitive, and that
\[
 \forall u\in a\;\forall v\in a\quad
       (u\Eq v\lor u\in v\lor v\in u).
\]
Here and below bounded quantifiers abbreviate ordinary first-order
formulas. When the members are saturated, this is a strict linear order
on their $\Eq$-classes. Define $\operatorname{EP}(s,t)$ to say that the
intersection of their member profiles, when nonempty, has endpoints:
\[
\begin{split}
 (\exists z\,(z\in s\land z\in t))\ \to\
 \exists l\exists g\,[&
 l\in s\land l\in t\land g\in s\land g\in t\\
 &\land\forall z\,((z\in s\land z\in t)\to
       ((l\Eq z\lor l\in z)\land(z\Eq g\lor z\in g)))].
\end{split}
\]
No intersection object is asserted to exist. The numerical predicate is
the following single formula:
\begin{equation}\label{un:number-definition}
\begin{split}
 \N(a):\Longleftrightarrow{}&
 \Sat(a)\land\Tr(a)\land\operatorname{Lin}(a)
 \land\forall u\,(u\in a\to\Sat(u))\\
 &\land\forall s\forall t\,
       [(t\Eq a\lor t\in a)\to\operatorname{EP}(s,t)].
\end{split}
\end{equation}
Its quantifiers range over existing objects, including objects in
nonstandard models; it asserts neither external finiteness nor external
well-foundedness. Finally put
\[
 \Emp(a):=\forall z\,z\notin a,\qquad
 \Succ(a,b):=\forall z\,[z\in b\leftrightarrow
                             (z\in a\lor z\Eq a)].
\]
Thus a successor, when it exists, adds the entire coextensionality class
of its predecessor.

\begin{lemma}[Local facts without successor totality]\label{un:local}
Empty objects exist and satisfy $\N$. The predicates $\Sat$ and $\N$
are invariant under $\Eq$. For numerical $a$:
\begin{enumerate}
\item $a\notin a$, and every member of $a$ is numerical;
\item if $\Succ(a,b)$ holds, then $b$ is numerical and $a$ is its
      greatest member class;
\item if $a$ is nonempty, it has a greatest member $d$, unique modulo
      $\Eq$, and $\Succ(d,a)$; its least members are empty.
\end{enumerate}
The relation $\Succ$ is invariant under $\Eq$ in its two arguments.
On numerical classes its existing outputs are unique, its existing
predecessors are unique, and no existing successor is empty.
None of these assertions presupposes that every number has a successor.
\end{lemma}
\begin{proof}
The unary false-matrix instance of $\UU$ has an output with no members.
An empty object satisfies every clause of~\eqref{un:number-definition}.
Saturation and numericality depend only on the member profile of their
argument, so are $\Eq$-invariant.

If $a\in a$, irreflexivity in $\operatorname{Lin}(a)$ applies to $a$
itself, giving a contradiction. Let $i\in a$. Then $\Sat(i)$ is part
of $\N(a)$, and transitivity gives $i\subseteq a$. If $u\in v\in i$,
all three of $u,v,i$ belong to $a$; transitivity of the restricted order
on $a$ gives $u\in i$. Hence $i$ is transitive. Its linear order and
saturation of its members are inherited. For an endpoint cutoff $t$,
if $t\in i$, transitivity puts $t$ in $a$; if $t\Eq i$, saturation of
$a$ also puts $t$ in $a$. The endpoint clause for $a$ therefore proves
the required clause for $i$.

Suppose $\Succ(a,b)$. Its member profile consists of the old members of
$a$ and one new $\Eq$-class, that of $a$. This profile is saturated.
Every old member is saturated, and any new member is coextensional with
the saturated object $a$. The profile is transitive: a member of an old
member lies in $a$, and a member of an object coextensional with $a$
also lies in $a$. Every old member lies below every new representative.
There is no reverse edge, because an edge from a representative of
$[a]$ into an old member, followed by transitivity and saturation, would
force $a\in a$. There is no edge within the new class for the same
reason. Thus the old strict class order extends by its one greatest
class.

For the endpoint condition of $b$, a cutoff $t\in b$ is either an old
member or coextensional with $a$, so the condition for $a$ applies. If
$t\Eq b$, its intersection with any $s$ is the old intersection with
$a$, possibly augmented by representatives of $[a]$. If the old part
is nonempty, its least endpoint remains least, and its greatest remains
greatest unless a new representative is present. If the old part is
empty, any present new representative is both endpoints. This proves
$\N(b)$ and the greatest-class assertion.

For nonempty $a$, apply its endpoint clause with $s=t=a$ to obtain a
greatest member $d$. It is numerical by heredity. Greatestness gives
$z\in a\to(z\in d\lor z\Eq d)$. Conversely, transitivity gives
$z\in d\to z\in a$, and saturation of $a$ gives
$z\Eq d\to z\in a$. This is $\Succ(d,a)$. If a least member $l$
had a member $z$, then $z\in a$, and leastness would give $l\Eq z$
or $l\in z$, each contradicting the strict class order since $z\in l$.
Thus $l$ is empty.

The definition of $\Succ$ proves its invariance and uniqueness of output
profiles. Each numerical predecessor is the greatest member class of
its output, so predecessors are unique as well. Every successor contains
its predecessor and is therefore nonempty.
\end{proof}

\begin{theorem}[Full representative induction before total successors]
\label{un:induction}
For every fixed ambient formula $\theta(v,\bar p)$, the theory $\UU$
proves the following induction implication, with arbitrary parameters
$\bar p$:
\begin{equation}\label{un:induction-form}
\begin{split}
 \bigl[&\forall e\,(\Emp(e)\to\theta(e,\bar p))\\
 &\land\forall a\forall b\,
   ((\N(a)\land\Succ(a,b)\land\theta(a,\bar p))
                            \to\theta(b,\bar p))\bigr]\\
 &\hspace{35mm}\to\forall n\,(\N(n)\to\theta(n,\bar p)).
\end{split}
\end{equation}
No invariance of $\theta$ under $\Eq$ is required.
\end{theorem}
\begin{proof}
First let $Q(v)$ have no other free variables. Assume its base and step
premises, and suppose $\N(n)\land\neg Q(n)$. Apply the legal unary
instance with matrix $\neg Q(x)$ and source $n$, obtaining $B$. For
every $x\in n$, transitivity gives $x\subseteq n$, so the two weak
filter clauses imply
\begin{equation}\label{un:local-failures}
               x\in B\quad\longleftrightarrow\quad\neg Q(x)
                  \qquad(x\in n).
\end{equation}

If the intersection of $B$ with $n$ is nonempty, take a least member
$i$. It is numerical by heredity and fails $Q$ by the upper filter
clause. It is not empty, by the base premise. Let $d$ be its greatest
member; Lemma~\ref{un:local} gives $\N(d)$ and $\Succ(d,i)$.
Transitivity puts $d$ in $n$. If $\neg Q(d)$, then
\eqref{un:local-failures} puts $d$ in the same intersection, contrary
to its least class $[i]$, since $d\in i$. Hence $Q(d)$, and the step
premise gives the contradiction $Q(i)$.

If that intersection is empty, every member of $n$ satisfies $Q$.
The counterexample $n$ is nonempty, so its greatest member $d$ satisfies
$Q(d)$ and $\Succ(d,n)$. The step premise again contradicts $\neg Q(n)$.
This second case is the crucial change: it uses $n$ itself as the least
failure and does not require an already existing successor of $n$.

For a formula with parameters, let $A_\theta(\bar p)$ be the conjunction
of its base and step premises in~\eqref{un:induction-form}. Form the
unary formula
\[
 Q(v):=\forall\bar p\,[A_\theta(\bar p)\to\theta(v,\bar p)],
\]
freshening every bound variable. Its base follows from the base component
of $A_\theta$. Given an existing successor and $Q$ at its predecessor,
fix any parameters satisfying $A_\theta$; the step component then gives
$Q$ at the successor. The unary result proves $Q$ for every number.
Specializing the bound parameters gives~\eqref{un:induction-form}.
The filter matrix used in this proof still has only its tested object
free.
\end{proof}

\section{Global least failures remove the source parameter}
\label{un:minimization}

For numerical arguments write
\[
                 a\preccurlyeq b\quad:\Longleftrightarrow\quad
                            (a\Eq b\lor a\in b).
\]

\begin{lemma}[Global comparison without successor totality]
\label{un:comparison}
For any two numbers $a,b$, exactly one of
$a\Eq b$, $a\in b$, and $b\in a$ holds. This conclusion can be
proved before totality of $\Succ$.
\end{lemma}
\begin{proof}
First prove by Theorem~\ref{un:induction}, on $b$, the conditional
no-gap property
\begin{equation}\label{un:no-gap}
 \N(a)\land\N(a')\land a\in b\land\Succ(a,a')
       \quad\to\quad(a'\in b\lor a'\Eq b).
\end{equation}
The empty case is vacuous. At an existing successor $\Succ(b,b')$,
membership $a\in b'$ means $a\in b$ or $a\Eq b$. In the first case,
the induction hypothesis puts $a'$ in $b$ or in $[b]$, so in $b'$.
In the second case, successor congruence gives $a'\Eq b'$.
This uses existing successors only.

Now induct on $a$ in the assertion that it compares with every number
$b$. An empty $a$ is coextensional with any empty $b$. If $b$ is
nonempty, its least member is empty; saturation of $b$ therefore gives
$a\in b$. For the step $\Succ(a,a')$, compare $a$ with $b$ by the
induction hypothesis. If $b\in a$ or $b\Eq a$, then $b\in a'$.
If $a\in b$, apply~\eqref{un:no-gap} to obtain $a'\in b$ or
$a'\Eq b$. This proves exhaustiveness. Irreflexivity, transitivity,
and saturation on numbers make the alternatives disjoint. Consequently
the numerical classes are strictly linearly ordered as claimed.
\end{proof}

\begin{lemma}[Global least numerical tuples]\label{un:least}
Every fixed formula $P(v,\bar q)$ invariant under $\Eq$ in its numerical
argument has a least numerical class whenever it has a numerical
witness. The parameters $\bar q$ may be arbitrary objects.
More generally, for each fixed positive integer $k$, every formula
$B(n_1,\ldots,n_k)$ invariant in its numerical arguments has a
lexicographically least numerical tuple whenever it has one. Its least
tuples are unique coordinatewise modulo $\Eq$ and are described by a
fixed first-order formula. If $B$ has no additional free parameter,
neither does this least-tuple formula.
\end{lemma}
\begin{proof}
For the first claim, apply full representative induction to the assertion
that the numerical class segment
$I_n=\{i:\N(i)\land i\preccurlyeq n\}$ either has no $P$-member or
has a least one. At an empty $n$, the segment consists of its one class,
and invariance of $P$ handles all representatives. At an existing
successor $\Succ(n,n')$, the definition gives
$I_{n'}=I_n\cup[n']$. If an old least witness exists it remains least;
otherwise the new class either supplies the least witness or supplies
none. This proves the induction assertion without totality of successors.
If $P$ has witness $n$, its least witness in $I_n$ is globally least:
comparison shows every number outside $I_n$ is larger than $n$.

For tuples, minimize the first coordinate for which an extension to a
$B$-tuple exists. Keep a chosen representative as a parameter of the
derived least-number theorem, and minimize the second coordinate among
its extensions; continue a fixed $k$ times. Invariance transports
extensions between equivalent choices of previous coordinates. The
result is a lexicographically least tuple and is unique coordinatewise
modulo $\Eq$.

For precision, let $\bar r<_{\mathrm{lex}}\bar s$ abbreviate the finite
disjunction
\[
 \bigvee_{j=1}^{k}
       \left[\left(\bigwedge_{i<j}r_i\Eq s_i\right)\land r_j\in s_j\right].
\]
One may take
\begin{equation}\label{un:least-formula}
\begin{split}
 L_B(\bar r):\Longleftrightarrow{}&
 \bigwedge_i\N(r_i)\land B(\bar r)\\
 &\land\forall\bar s\,
 \left[\left(\bigwedge_i\N(s_i)\land B(\bar s)\right)\to
   \left(\bigwedge_i r_i\Eq s_i\ \lor\
                             \bar r<_{\mathrm{lex}}\bar s\right)\right].
\end{split}
\end{equation}
Tuple notation here abbreviates a fixed finite list of variables, not
objects coding tuples. Formula~\eqref{un:least-formula} and its proof
are obtained uniformly from the syntax of $B$.
\end{proof}

\begin{theorem}[Bounded numerical comprehension]\label{un:bnc}
Let $\chi(z,\bar n)$ be any fixed formula whose only free variables are
the displayed ones, with $k$ numerical parameter positions. Let
$I_\chi$ assert invariance of $\chi$ when those parameters are replaced
by coextensional numerical representatives, and let
\[
 S_\chi:\quad
 \forall\bar n\in\N\;\exists a\;\forall z\,
                       (\chi(z,\bar n)\to z\subseteq a).
\]
Then $\UU$ proves the conditional theorem
\begin{equation}\label{un:bnc-statement}
 (I_\chi\land S_\chi)\ \to\
 \forall\bar n\in\N\;\exists b\;\forall z\,
                       (z\in b\leftrightarrow\chi(z,\bar n)).
\end{equation}
No invariance condition is imposed on the tested object $z$.
\end{theorem}
\begin{proof}
Assume $I_\chi$ and $S_\chi$. For $k>0$, suppose an output fails to
exist at some numerical tuple. Define
\[
 B(\bar n):=\neg\exists b\;\forall z\,
                          (z\in b\leftrightarrow\chi(z,\bar n)).
\]
The assumed invariance of $\chi$ makes $B$ invariant in its numerical
arguments. Lemma~\ref{un:least} supplies a globally least failing tuple
$\bar r$. Use the following legal unary matrix:
\begin{equation}\label{un:bnc-matrix}
           \varphi(z):=\exists\bar n\,
                         [L_B(\bar n)\land\chi(z,\bar n)].
\end{equation}
Every tuple satisfying $L_B$ is coordinatewise coextensional with
$\bar r$, so $I_\chi$ proves
$\varphi(z)\leftrightarrow\chi(z,\bar r)$ for every ambient object
$z$. Choose a support source $a$ from $S_\chi$ at $\bar r$. Apply
the unary weak filter to~\eqref{un:bnc-matrix} at that source. Its
upper clause gives $z\in b\to\chi(z,\bar r)$. Conversely,
$\chi(z,\bar r)$ implies both $z\subseteq a$ and $\varphi(z)$, so
the lower clause gives $z\in b$. This is the forbidden exact output.

For $k=0$, the formula $\chi(z)$ is already unary. Choose any supporting
source from $S_\chi$ and apply its weak-filter instance directly; the
same two implications give exactness.
\end{proof}

The support source is an existentially chosen object in the proof; it
does not occur free in~\eqref{un:bnc-matrix}. The argument proves a
derived conditional comprehension theorem. It neither exactifies all
source-free instances nor licenses arbitrary ambient parameters. A use
of the theorem must supply proofs of its stated invariance and support
premises. The least-failure predicate is a finite formula formed from
$\chi$; no algorithm evaluates it, and no self-referential truth
predicate is introduced.

\begin{corollary}[The bootstrap is complete]\label{un:bootstrap}
Every number has a numerical successor. Every number $n$ has an exact
power witness
\[
             \Pow(n,W):=\forall z\,(z\in W\leftrightarrow z\subseteq n).
\]
Every finite tuple of numbers is contained in a common numerical object.
\end{corollary}
\begin{proof}
Apply Theorem~\ref{un:bnc} first to
$\chi(z,n):=(z\in n\lor z\Eq n)$. It is invariant in $n$.
Every selected object is a subset of $n$, by transitivity or
coextensionality, so $n$ itself is a support source. Its exact output
satisfies $\Succ(n,b)$ and is numerical by Lemma~\ref{un:local}.

Next use $\chi(z,n):=(z\subseteq n)$, again invariant in $n$ and
supported at $n$. The output is the asserted exact power. Only powers
of numerical sources have been proved here; no power axiom for all
ambient objects is being assumed.

Finally, numerical comparison chooses a greatest class from any fixed
finite tuple. A successor of its representative contains every entry
of that tuple, by saturation and the successor profile.
\end{proof}

The order of these arguments rules out the tempting circular proof:
local predecessor facts imply induction; induction implies comparison
and least failures; least failures imply bounded numerical comprehension;
only then does comprehension imply total successors and numerical
powers. Neither a preliminary extensional core nor a source-recovery
decoder is needed.

